%%%%%
%%%%%  Naudokite LUALATEX, ne LATEX.
%%%%%
%%%%
\documentclass[]{VUMIFTemplateClass}

\usepackage{indentfirst}
\usepackage{amsmath, amsthm, amssymb, amsfonts}
\usepackage{mathtools}
\usepackage{physics}
\usepackage{graphicx}
\usepackage{verbatim}
\usepackage[hidelinks]{hyperref}
\usepackage{color,algorithm,algorithmic}
\usepackage[nottoc]{tocbibind}
\usepackage{tocloft}
\usepackage[version=4]{mhchem}
\usepackage{textgreek}
\usepackage{textcomp}
\usepackage{float}
\usepackage{booktabs}
\usepackage{listings}
\lstset{
    language=Python,
    basicstyle=\ttfamily\footnotesize,
    keywordstyle=\bfseries,
    commentstyle=\itshape,
    stringstyle=\ttfamily,
    breaklines=true,
    frame=single,
    showstringspaces=false,
    numbers=left,
    numberstyle=\tiny,
    xleftmargin=1.5em,
}

\usepackage{titlesec}
\newcommand{\sectionbreak}{\clearpage}

\makeatletter
\renewcommand{\fnum@algorithm}{\thealgorithm}
\makeatother
\renewcommand\thealgorithm{\arabic{algorithm} algoritmas}

\usepackage{biblatex}
\bibliography{bibliografija}

\newcommand{\EE}{\mathbb{E}\,}
\newcommand{\ee}{{\mathrm e}}
\newcommand{\RR}{\mathbb{R}}

\studijuprograma{Informatikos}
\darbotipas{Ataskaita}
\darbopavadinimas{Kietųjų dalelių koncentracijos prognozavimas laiko eilutėse}
\darbopavadinimasantras{Duomenų paruošimas}
\autorius{Mindaugas Žiemys}

\vadovas{Linas Litvinas, Asist., Dr.}

\begin{document}
    \selectlanguage{lithuanian}

    \onehalfspacing
    \input{titulinis}

    \singlespacing

    \tableofcontents
    \onehalfspacing


    \section{Duomenų paruošimas}
    Pradinis duomenų rinkinys apima 2012--2025 metų valandinius oro taršos matavimus iš Londono Merilebono kelio stoties -- iš viso 122\,736 eilutės.
    Duomenų paruošimo seką apima šie pagrindiniai žingsniai:
    \begin{itemize}
        \item \ce{CO} stulpelio pašalinimas.
        \item COVID--19 pandemijos 2020--2022 m. laikotarpio pašalinimas.
        \item Neigiamų reikšmių pašalinimas.
        \item Eilučių su trūkstamomis reikšmėmis pašalinimas.
        \item Vėjo krypties $\sin$ ir $\cos$ dedamųjų pridėjimas.
        \item Dienos kategorijos stulpelio sukūrimas.
        \item Min-Max normalizavimas -- verčių perkėlimas į intervalą $[-1, 1]$.
    \end{itemize}

    \subsection{CO pašalinimas}
    Iš visų 122\,736 eilučių, 30\,770 (25,1\,\%) turi bent vieną trūkstamą reikšmę.
    \ref{tab:missing} lentelėje pateikiamas eilučių su trūkstamomis reikšmėmis kiekis ir jų dalis, skaičiuojama nuo viso pradinio duomenų rinkinio ($N=122\,736$).

    \begin{table}[H]
        \centering
        \caption{Trūkstamos reikšmės pagal atributą ($N = 122\,736$)}
        \small
        \setlength{\tabcolsep}{6pt}
        \begin{tabular}{|l|r|r|}
            \hline
            \textbf{Atributas} & \textbf{Kiekis} & \textbf{Dalis (\%)} \\
            \hline
            \ce{CO}            & 12\,704         & 10,35               \\
            \ce{KD_{10}}       & 10\,379         & 8,46                \\
            \ce{KD_{2,5}}      & 9\,649          & 7,86                \\
            \ce{O3}            & 4\,359          & 3,55                \\
            \ce{NO}            & 3\,653          & 2,98                \\
            \ce{NO2}           & 3\,653          & 2,98                \\
            Vėjo greitis       & 3\,168          & 2,58                \\
            Temperatūra        & 3\,168          & 2,58                \\
            Vėjo kryptis       & 3\,168          & 2,58                \\
            \hline
        \end{tabular}
        \label{tab:missing}
    \end{table}    \ce{CO} turi didžiausią trūkstamų reikšmių dalį (10,35\,\%). Dėl šių priežasčių \ce{CO} pašalintas iš duomenų rinkinio.
    Pašalinus \ce{CO}, eilučių su trūkstamomis reikšmėmis liko 21\,636 (17,6\,\%).

    \subsection{COVID--19 pandemijos 2020--2022 m. laikotarpio pašalinimas}
    Sprendimą pašalinti šio laikotarpio duomenis lėmė dvi priežastys.
    Pirma, COVID--19 pandemijos metu dėl taikytų ribojimų oro taršos tendencijos buvo netipiškos.
    Antra, šiuo laikotarpiu fiksuotas žymiai didesnis \ce{KD_{2,5}} duomenų trūkumas (žr. \ref{tab:pm25_yearly_comparison} lentelę).
    \begin{table}[H]
        \centering
        \caption{Vidutinis trūkstamų \ce{KD_{2,5}} reikšmių kiekis ir dalis per metus}
        \small
        \setlength{\tabcolsep}{12pt}
        \begin{tabular}{|l|c|r|}
            \hline
            \textbf{Laikotarpis}                & \textbf{Kiekis (vid.)} & \textbf{Dalis (\%)} \\
            \hline
            2012--2019 m.                       & 483                    & 5,51                \\
            \textbf{2020--2022 m. (Pandeminis)} & \textbf{1\,476}        & \textbf{16,84}      \\
            2023--2025 m.                       & 436                    & 4,97                \\
            \hline
        \end{tabular}
        \label{tab:pm25_yearly_comparison}
    \end{table}
    Pašalinus COVID--19 laikotarpio įrašus, pilnas duomenų rinkinį sudaro 96\,432 įrašai, iš jų 13\,875 (14,4\%) -- turi bent vieną trūkstamą reikšmę.

    \subsection{Neigiamų reikšmių pašalinimas}
    Kadangi koncentracija fiziškai negali būti mažesnė už nulį, tokie duomenys laikomi matavimo įrangos paklaidomis ir yra pašalinami.
    \begin{table}[H]
        \centering
        \caption{Teršalų koncentracijų ribinės reikšmės ir neigiamų įrašų kiekis}
        \small
        \setlength{\tabcolsep}{8pt}
        \begin{tabular}{|l|r|r|c|}
            \hline
            \textbf{Atributas} & \textbf{Min} & \textbf{Max} & \textbf{Neigiamų kiekis} \\
            \hline
            \ce{O3}            & $-0,93$      & 156,46       & 22                       \\
            \ce{NO}            & $-0,37$      & 872,83       & 9                        \\
            \ce{PM_{10}}       & $-2,90$      & 187,90       & 44                       \\
            \ce{PM_{2,5}}      & $-5,00$      & 127,60       & 280                      \\
            \ce{NO2}           & 0,00         & 321,91       & 0                        \\
            \hline
            \textbf{Iš viso}   & --           & --           & \textbf{355}             \\
            \hline
        \end{tabular}
        \label{tab:negative_concentrations}
    \end{table}
    Pašalinus neigiamas reikšmes, duomenų rinkinį sudaro 96\,211 įrašai, iš jų 13\,855 (14,4\%) -- turi bent vieną trūkstamą reikšmę.

    \subsection{Eilučių su trūkstamomis reikšmėmis pašalinimas}
    Analizė parodo, kad 69,1\% visų duomenų spragų (žr. \ref{tab:gap_distribution} lentelę) sudaro laikotarpiai, ilgesni nei 24 val.

    \begin{table}[H]
        \centering
        \caption{Įrašų su trūkstamomis reikšmėmis pasiskirstymas pagal trukmę}
        \small
        \setlength{\tabcolsep}{8pt}
        \begin{tabular}{|l|r|r|}
            \hline
            \textbf{Tarpo ilgis} & \textbf{Eilučių kiekis} & \textbf{Dalis (\%)} \\
            \hline
            1 val.               & 470                     & 3,4                 \\
            2 val.               & 301                     & 2,2                 \\
            3--5 val.            & 322                     & 2,3                 \\
            6--12 val.           & 573                     & 4,1                 \\
            12--24 val.          & 2\,612                  & 18,9                \\
            24+ val.             & 9\,577                  & 69,1                \\
            \hline
            \textbf{Iš viso}     & \textbf{13\,855}        & \textbf{100,0}      \\
            \hline
        \end{tabular}
        \label{tab:gap_distribution}
    \end{table}
    Visos eilutės su trūkstamomis reikšmėmis yra pašalinamos.
    Pašalinus trūkstamus įrašus, pilną duomenų rinkinį sudaro 82\,356 įrašai.

    \subsection{Vėjo krypties sin ir cos dedamųjų pridėjimas.}

    Vėjo kryptis (laipsniais, 0--360\textdegree) konvertuota į dvi dedamąsias: sinusą ir kosinusą:

    \begin{equation}
        s = \sin\!\left(\alpha \cdot \frac{\pi}{180}\right), \qquad
        c = \cos\!\left(\alpha \cdot \frac{\pi}{180}\right)
        \label{eq:windsincos}
    \end{equation}

    Trigonometrinės funkcijos $\sin$ ir $\cos$ naudojamos su radianų argumentais, todėl vėjo krypties matavimai laipsniais ($\alpha$) konvertuojami taikant $\pi/180$.

    Transformavimas į $\sin$ ir $\cos$ išsprendžia cikliškumo problemą.
    Kampai 0\textdegree{} ir 360\textdegree{} yra identiški, tačiau skaitiškai skiriasi: $|360 - 0| = 360$, nors faktinis skirtumas 0\textdegree.
    Sinuso ir kosinuso atvaizdavime šie kampai duoda identiškas reikšmių poras.
    Originalus vėjo krypties stulpelis pašalintas.

    \subsection{Dienos kategorijos sukūrimas}
    Kiekviena laiko eilutė susieta su dienos kategorija, kurią galima aprašyti:

    \begin{equation}
        \text{kategorija}(d) =
        \begin{cases}
            1, & \text{jei } d \text{ yra poilsio diena} \\
            2, & \text{jei } d+1 \text{ yra poilsio diena} \\
            0, & \text{kitais atvejais (darbo diena)}
        \end{cases}
        \label{eq:daycat}
    \end{equation}

    čia poilsio diena -- savaitgalis (šeštadienis arba sekmadienis) arba JK valstybinė šventė.
    Dienos kategorija bus naudojamas kaip filtras -- prognozei naudojami įrašai apribojami tik tais įrašais, kurie priklauso tai pačiai kategorijai.
    Tokiu būdu užtikrinama, kad, pavyzdžiui, sekmadienio tarša nebūtų prognozuojama remiantis pirmadienio duomenimis, nes transporto intensyvumas šiomis dienomis skiriasi.

    \subsection{Sezoniniai teršalų trendai}

    \ref{tab:seasonal_trends} lentelėje pateikiami pagrindinių teršalų mėnesiniai vidurkiai.
    Sausio -- kovo mėn. vidutinė \ce{KD_{2,5}} koncentracija išauga. NO koncentracija padidėja lapkričio -- sausio mėn.

    \begin{table}[H]
        \centering
        \caption{Teršalų vidutinės mėnesinės koncentracijos (\textmu g/m\textsuperscript{3})}
        \small
        \setlength{\tabcolsep}{5pt}
        \begin{tabular}{|l|r|r|r|r|r|}
            \hline
            \textbf{Mėnuo} & \textbf{\ce{O3}} & \textbf{\ce{NO}} & \textbf{\ce{NO2}} & \textbf{\ce{KD_{10}}} & \textbf{\ce{KD_{2,5}}} \\
            \hline
            Sausis         & 15,25            & \textbf{122,16}  & 76,36             & 26,67                 & \textbf{17,08}         \\
            Vasaris        & 18,50            & 108,38           & 72,93             & 27,38                 & \textbf{17,82}         \\
            Kovas          & 23,71            & 92,03            & 73,72             & 29,23                 & \textbf{19,99}         \\
            Balandis       & 30,55            & 81,20            & 72,36             & 25,27                 & 16,63                  \\
            Gegužė         & 32,42            & 75,45            & 68,53             & 22,69                 & 15,11                  \\
            Birželis       & 27,13            & 87,40            & 72,36             & 20,49                 & 13,16                  \\
            Liepa          & 22,48            & 88,32            & 72,83             & 19,62                 & 12,91                  \\
            Rugpjūtis      & 20,33            & 83,07            & 68,15             & 21,05                 & 12,94                  \\
            Rugsėjis       & 18,00            & 95,17            & 70,47             & 22,46                 & 14,94                  \\
            Spalis         & 14,60            & 108,68           & 71,47             & 23,83                 & 14,58                  \\
            Lapkritis      & 13,09            & \textbf{121,73}  & 71,55             & 25,44                 & 15,77                  \\
            Gruodis        & 16,03            & \textbf{122,40}  & 73,90             & 24,31                 & 15,03                  \\
            \hline
        \end{tabular}
        \label{tab:seasonal_trends}
    \end{table}
    Sezoniniai skirtumai rodo, kad modelio prognozavimui įtaką gali turėti sezoniškumas.
    Sezoniškumas bus įtrauktas į K-NN metodą, apribojant paiešką $\pm 20$ dienų tolerancijos intervalu aplink tikslinę datą, per visų turimų metų duomenų rinkinį.
    Pavyzdžiui, prognozuojant 2023 m. rugpjūčio 2 d. taršą, paieškoa apima visų metų laikotarpius nuo liepos 13 d. iki rugpjūčio 22 d.

    \subsection{Min-max normalizavimas}

    Skaitiniai matavimų atributai normalizuoti į $[-1, 1]$ intervalą pagal formulę:

    \begin{equation}
        x' = \frac{2(x - x_{\min})}{x_{\max} - x_{\min}} - 1
        \label{eq:minmax}
    \end{equation}

    čia $x$ -- originali reikšmė, $x_{\min}$ ir $x_{\max}$ -- atitinkamo atributo minimali ir maksimali reikšmės visoje duomenų aibėje, $x'$ -- normalizuota reikšmė.
    Normalizavimas atliekamas kiekvienam atributui atskirai, naudojant tik jam būdingas minimalias ir maksimalias reikšmes.
    Tokiu būdu visų kintamųjų vertės suvienodinamas iki $[-1, 1]$ intervalo.

    \ref{tab:minmax_final} lentelėje pateikiamos normalizuojamų atributų minimalios ir maksimalios reikšmės.
    Neigiamos minimalios reikšmės buvo pašalintos praeitame žingsnyje.

    \begin{table}[H]
        \centering
        \caption{Atributų minimalios ir maksimalios reikšmės, naudojamos normalizavimui}
        \small
        \setlength{\tabcolsep}{8pt}
        \begin{tabular}{|l|c|r|r|}
            \hline
            \textbf{Atributas} & \textbf{Vnt.}                  & \textbf{$x_{\min}$} & \textbf{$x_{\max}$} \\
            \hline
            \ce{O3}            & \textmu g/m\textsuperscript{3} & 0,00                & 151,47              \\
            \ce{NO}            & \textmu g/m\textsuperscript{3} & 0,00                & 872,83              \\
            \ce{NO2}           & \textmu g/m\textsuperscript{3} & 0,00                & 321,91              \\
            \ce{KD_{10}}       & \textmu g/m\textsuperscript{3} & 0,00                & 187,90              \\
            \ce{KD_{2,5}}      & \textmu g/m\textsuperscript{3} & 0,00                & 127,60              \\
            Vėjo greitis       & m/s                            & 0,00                & 13,70               \\
            Temperatūra        & \textdegree C                  & $-10,50$            & 33,40               \\
            \hline
        \end{tabular}
        \label{tab:minmax_final}
    \end{table}
    Vėjo krypties sinuso ir kosinuso komponentai nebuvo normalizuoti papildomai, nes pagal apibrėžimą jau patenka į $[-1,\,1]$ intervalą.
    Dienos kategorija ir valanda yra kategoriniai atributai, jų normalizuoti nereikia.


    \section{Prognozavimo metodika}

    Oro taršos prognozavimui naudojamas modifikuotas $K$ artimiausių kaimynų (KNN) metodas su svoriais.
    Modelio esmė -- rasti istorinius 5~valandų laikotarpius, kurių meteorologinės ir taršos sąlygos yra panašiausios į dabartinę situaciją, ir pagal jų tęsinį prognozuoti \ce{KD_{2,5}} koncentraciją ateinančioms 12~valandų.

    Panašumui įvertinti naudojamas svertinis Euklido atstumas, sudarytas iš penkių atskirų atstumo funkcijų.
    Kiekviena funkcija apima tam tikrą atributų poaibį ir matuoja atstumą tarp dabartinio 5~valandų lango $\mathbf{x}$ ir istorinio lango $\mathbf{y}$.
    Funkcijoms priskiriami svorio koeficientai, o galutinis panašumas skaičiuojamas kaip svertinė suma.

    \subsection{Atstumo funkcijos}

    \subsubsection*{1 funkcija. Vėjo dinamika ($d_1$)}
    Apima vėjo greitį ($v$) ir vėjo krypties dedamąsias ($s = \sin\alpha$, $c = \cos\alpha$).
    Vėjas yra pagrindinis teršalų sklaidos horizontalioje kryptyje veiksnys.
    Iš viso Euklido atstumas skaičiuojamas per 3 \times 5 = 15 matmenų.
    \begin{equation}
        d_1 = \sqrt{\sum_{h=1}^{5} \left[ (s_{x,h} - s_{y,h})^2 + (c_{x,h} - c_{y,h})^2 + (v_{x,h} - v_{y,h})^2 \right]}
    \end{equation}

    \subsubsection*{2 funkcija. Temperatūra ($d_2$)}
    Temperatūra veikia teršalų sklaidą vertikalioje kryptyje.
    Apima 5 temperatūros ($T$) matavimus.
    \begin{equation}
        d_2 = \sqrt{\sum_{h=1}^{5} (T_{x,h} - T_{y,h})^2}
    \end{equation}

    \subsubsection*{3 funkcija. Dujiniai teršalai ($d_3$)}
    Apima visus dujinius teršalus, išskyrus kietąsias daleles: \ce{O3}, \ce{NO} ir \ce{NO2} koncentracijas.
    Iš viso: $3 \text{ atributai} \times 5 \text{ valandos} = 15$ reikšmių.
    \begin{equation}
        d_3 = \sqrt{\sum_{h=1}^{5} \left[ ({\ce{O3}}_{x,h} - {\ce{O3}}_{y,h})^2 + ({\ce{NO}}_{x,h} - {\ce{NO}}_{y,h})^2 + ({\ce{NO2}}_{x,h} - {\ce{NO2}}_{y,h})^2 \right]}
    \end{equation}

    \subsubsection*{4 funkcija. \ce{KD_{10}} ($d_4$)}
    Apima 5 stambiųjų kietųjų dalelių \ce{KD_{10}} matavimus.
    \begin{equation}
        d_4 = \sqrt{\sum_{h=1}^{5} ({\ce{KD_{10}}}_{x,h} - {\ce{KD_{10}}}_{y,h})^2}
    \end{equation}

    \subsubsection*{5 funkcija. \ce{KD_{2,5}} dabartinis lygis ($d_5$)}
    Apima 5 tikslinio kintamojo \ce{KD_{2,5}} matavimus.
    Užtikrina, kad rastas kaimynas turėtų panašų dabartinį \ce{KD_{2,5}} lygį.
    \begin{equation}
        d_5 = \sqrt{\sum_{h=1}^{5} ({\ce{KD_{2,5}}}_{x,h} - {\ce{KD_{2,5}}}_{y,h})^2}
    \end{equation}

    \subsection{Bendra atstumo funkcijos išraiška}

    Visos penkios funkcijos turi tą pačią struktūrą -- Euklido atstumą per 5 valandų langą.
    Bet kurią iš jų galima užrašyti bendra forma:
    \begin{equation}
        d_j = \sqrt{\sum_{h=1}^{5} \sum_{f \in F_j} (x_{h,f} - y_{h,f})^2}
        \label{eq:dist_group}
    \end{equation}
    čia $F_j$ -- $j$-osios funkcijos atributų aibė (pvz., $F_1 = \{v, s, c\}$, $F_2 = \{T\}$), $x_{h,f}$ ir $y_{h,f}$ -- atitinkamai dabartinė ir istorinė normalizuoto atributo $f$ reikšmė $h$~valandą.

    \subsection{Agreguotas svertinis atstumas}

    Visos penkios funkcijos apjungiamos į vieną svertinę sumą:
    \begin{equation}
        D = w_1 d_1 + w_2 d_2 + w_3 d_3 + w_4 d_4 + w_5 d_5
        \label{eq:total_dist}
    \end{equation}
    čia $w_j$ -- svorio koeficientai, tenkinantys sąlygą $\sum_{j=1}^{5} w_j = 1$.
    Svoriai leidžia reguliuoti kiekvienos funkcijos įtaką bendram panašumui.
    Pradiniuose eksperimentuose visiems svoriams bus priskirta vienoda reikšmė $w_j = 0{,}2$.
    Optimalūs svoriai bus parinkti eksperimentų metu, siekiant minimizuoti prognozės paklaidą.

    \subsection{Paieškos erdvės apribojimai}

    Potencialių kaimynų paieška istoriniuose duomenyse apribojama šiais filtrais:
    \begin{itemize}
        \item \textbf{Laiko sinchronizacija:} lyginami tik tie istoriniai langai, kurie prasideda tą pačią paros valandą kaip ir dabartinė situacija. Pavyzdžiui, jei dabartinis langas apima 8--12~val., istoriniai langai taip pat turi prasidėti nuo 8~val.
        \item \textbf{Dienos kategorija:} darbo dienų langai lyginami tik su darbo dienų langais, savaitgalių -- tik su savaitgalių.
        \item \textbf{Sezoniškumas:} paieška apribojama $\pm 20$~dienų intervalu aplink tikslinę datą per visus turimus metus. Pavyzdžiui, prognozuojant rugpjūčio 2~d. taršą, paieška apima laikotarpius nuo liepos 13~d. iki rugpjūčio 22~d. visais turimais metais.
        \item \textbf{Duomenų vientisumas:} įtraukiami tik tie blokai, kuriuose nėra trūkių nei 5~valandų įvesties, nei 12~valandų prognozės lange. Pavyzdžiui, jei istorinis langas prasideda 8~val., tai visos įvesties reikšmės nuo 8 iki 12~val. ir prognozės nuo 13 iki 24~val. turi būti be trūkstamų reikšmių.
    \end{itemize}

    \subsection{Slenksčio taikymas ir baltosios dėžės principas}

    Modelis remiasi baltosios dėžės (\textit{White Box}) principu -- kaimynų atranka pagrįsta ne fiksuotu skaičiumi $K$, o slenksčio konstanta $S$:

    \begin{equation}
        \text{Sprendimas} =
        \begin{cases}
            \text{Prognozuoti}, & \text{jei } \min(D) \le S \\
            \text{Neprognozuoti}, & \text{jei } \min(D) > S
        \end{cases}
        \label{eq:threshold}
    \end{equation}

    Jei nė vienas istorinis langas netenkiną slenksčio sąlygos, modelis informuoja, kad prognozė negali būti sudaryta -- istorijoje nėra pakankamai panašios situacijos.
    Tai esminis skirtumas nuo juodosios dėžės metodų (pvz., dirbtinių neuroninių tinklų), kurie prognozę pateikia visada, neatsižvelgiant į tai, ar ji pagrįsta analogija su istoriniais duomenimis.
    Slenksčio konstanta $S$ taip pat bus parenkama eksperimentiškai -- ieškant balanso tarp prognozės tikslumo ir padengimo (dalies atvejų, kuriems modelis pateikia prognozę).

    \subsection{Prognozės sudarymas}

    Atrenkami visi kaimynai, tenkinantys sąlygą $D \le S$.
    Galutinė 12~valandų \ce{KD_{2,5}} prognozė kiekvienam laiko žingsniui $t \in \{1, \dots, 12\}$ skaičiuojama kaip atrinktų kaimynų ateities reikšmių vidurkis:

    \begin{equation}
        \hat{y}_{t} = \frac{1}{K} \sum_{k=1}^{K} y_{t}^{(k)}
        \label{eq:forecast}
    \end{equation}
    čia $K$ -- kaimynų, patenkančių į slenksčio zoną, skaičius, $y_{t}^{(k)}$ -- $k$-ojo kaimyno užfiksuota \ce{KD_{2,5}} koncentracija po $t$~valandų nuo atskaitos taško.


    \section{Rezultatai}

    \subsection{Validžių duomenų iškarpa}
    Lentelėje \ref{tab:raw_data_sample} pateikiama validžių originalių duomenų iškarpa prieš normalizavimą.
    \begin{table}[H]
        \centering
        \caption{Validžių originalių duomenų iškarpa}
        \small
        \resizebox{\textwidth}{!}{
            \begin{tabular}{|l|c|r|r|r|r|r|r|r|r|r|}
                \hline
                \textbf{Data} & \textbf{Laikas} & \textbf{\ce{O3}} & \textbf{\ce{NO}} & \textbf{\ce{NO2}} & \textbf{\ce{CO}} & \textbf{\ce{KD_{10}}} & \textbf{\ce{KD_{2,5}}} & \textbf{Vėjo krypt.} & \textbf{Vėjo greit.} & \textbf{Temp.} \\
                \hline
                2016-03-30    & 06:00:00        & 3,79             & 289,29           & 134,73            & 0,595            & 31,9                  & 14,17                  & 246,8                & 2,6                  & 1,6            \\
                2016-03-30    & 07:00:00        & 2,94             & 404,42           & 160,86            & 0,863            & 35,2                  & 16,35                  & 253,6                & 3,5                  & 3,4            \\
                2016-03-30    & 08:00:00        & 3,04             & 442,75           & 169,61            & 0,952            & 35,4                  & 16,95                  & 257,8                & 3,5                  & 6,2            \\
                2016-03-30    & 09:00:00        & 3,44             & 455,46           & 194,32            & 0,833            & 45,0                  & 18,23                  & 264,3                & 4,2                  & 8,5            \\
                2016-03-30    & 10:00:00        & 5,09             & 391,30           & 167,55            & 0,862            & 43,9                  & 16,95                  & 270,5                & 4,8                  & 10,0           \\
                2016-03-30    & 11:00:00        & 9,08             & 231,78           & 138,45            & 0,595            & 27,7                  & 11,00                  & 266,9                & 4,6                  & 11,0           \\
                2016-03-30    & 12:00:00        & 7,28             & 282,47           & 148,30            & 0,714            & 37,1                  & 16,25                  & 267,2                & 4,3                  & 11,8           \\
                2016-03-30    & 13:00:00        & 8,68             & 265,75           & 139,28            & 0,535            & 32,4                  & 16,25                  & 259,4                & 4,3                  & 12,0           \\
                2016-03-30    & 14:00:00        & 12,27            & 159,92           & 104,98            & 0,417            & 19,3                  & 8,92                   & 255,7                & 4,2                  & 11,7           \\
                2016-03-30    & 15:00:00        & 13,02            & 164,73           & 108,93            & 0,476            & 20,5                  & 10,21                  & 252,9                & 3,7                  & 11,1           \\
                2016-03-30    & 16:00:00        & 12,17            & 182,46           & 135,16            & 0,684            & 28,4                  & 10,60                  & 256,4                & 3,3                  & 10,4           \\
                2016-03-30    & 17:00:00        & 7,83             & 279,46           & 141,91            & 0,952            & 30,4                  & 13,58                  & 250,0                & 1,8                  & 9,2            \\
                2016-03-30    & 18:00:00        & 13,32            & 145,24           & 114,01            & 0,595            & 22,1                  & 12,49                  & 227,0                & 1,8                  & 7,5            \\
                \hline
            \end{tabular}
        }
        \label{tab:raw_data_sample}
    \end{table}

    Lentelėje \ref{tab:normalized_sample} pateikiama šios iškarpos normalizuota ir transformuota versija.
    \begin{table}[H]
        \centering
        \caption{Validžių normalizuotų ir transformuotų duomenų škarpa}
        \small
        \resizebox{\textwidth}{!}{
            \begin{tabular}{|l|c|r|r|r|r|r|r|r|c|r|r|}
                \hline
                \textbf{Data} & \textbf{Val.} & \textbf{\ce{O3}} & \textbf{\ce{NO}} & \textbf{\ce{NO2}} & \textbf{\ce{KD_{10}}} & \textbf{\ce{KD_{2,5}}} & \textbf{Vėjo gr.} & \textbf{Temp.} & \textbf{Dienos kat.} & \textbf{Vėjo $\sin$} & \textbf{Vėjo $\cos$} \\
                \hline
                2016-03-30    & 6             & -0,94            & -0,34            & -0,16             & -0,66                 & -0,78                  & -0,62             & -0,45          & 0                    & -0,92                & -0,39                \\
                2016-03-30    & 7             & -0,95            & -0,07            & 0,00              & -0,63                 & -0,74                  & -0,49             & -0,37          & 0                    & -0,96                & -0,28                \\
                2016-03-30    & 8             & -0,95            & 0,01             & 0,05              & -0,62                 & -0,73                  & -0,49             & -0,24          & 0                    & -0,98                & -0,21                \\
                2016-03-30    & 9             & -0,94            & 0,04             & 0,21              & -0,52                 & -0,71                  & -0,39             & -0,13          & 0                    & -1,00                & -0,10                \\
                2016-03-30    & 10            & -0,92            & -0,10            & 0,04              & -0,53                 & -0,73                  & -0,30             & -0,07          & 0                    & -1,00                & 0,01                 \\
                2016-03-30    & 11            & -0,85            & -0,47            & -0,14             & -0,71                 & -0,83                  & -0,33             & -0,02          & 0                    & -1,00                & -0,05                \\
                2016-03-30    & 12            & -0,88            & -0,35            & -0,08             & -0,61                 & -0,75                  & -0,37             & 0,02           & 0                    & -1,00                & -0,05                \\
                2016-03-30    & 13            & -0,86            & -0,39            & -0,13             & -0,66                 & -0,75                  & -0,37             & 0,03           & 0                    & -0,98                & -0,18                \\
                2016-03-30    & 14            & -0,80            & -0,63            & -0,35             & -0,79                 & -0,86                  & -0,39             & 0,01           & 0                    & -0,97                & -0,25                \\
                2016-03-30    & 15            & -0,79            & -0,62            & -0,32             & -0,78                 & -0,84                  & -0,46             & -0,02          & 0                    & -0,96                & -0,29                \\
                2016-03-30    & 16            & -0,80            & -0,58            & -0,16             & -0,70                 & -0,83                  & -0,52             & -0,05          & 0                    & -0,97                & -0,24                \\
                2016-03-30    & 17            & -0,87            & -0,36            & -0,12             & -0,68                 & -0,79                  & -0,74             & -0,10          & 0                    & -0,94                & -0,34                \\
                2016-03-30    & 18            & -0,78            & -0,67            & -0,29             & -0,76                 & -0,80                  & -0,74             & -0,18          & 0                    & -0,73                & -0,68                \\
                2016-03-30    & 19            & -0,68            & -0,87            & -0,46             & -0,87                 & -0,91                  & -0,75             & -0,24          & 0                    & -0,76                & -0,65                \\
                \hline
            \end{tabular}        }
        \label{tab:normalized_sample}
    \end{table}

    \subsection{Iškarpa su nevalidžiais duomenimis}

    Lentelėje \ref{tab:missing_data} pateikiama duomenų iškarpa su trūkstamomis reikšmėmis 11, 12, 13 valandomis bei su neigiama \ce{KD_{2,5}} reikšme 6 valandą.
    \begin{table}[H]
        \centering
        \caption{Originalių duomenų iškarpa su trūkstamomis ir neigiamomis reikšmėmis}
        \small
        \resizebox{\textwidth}{!}{
            \begin{tabular}{|l|c|r|r|r|r|r|r|r|r|r|}
                \hline
                \textbf{Data} & \textbf{Laikas} & \textbf{\ce{O3}} & \textbf{\ce{NO}} & \textbf{\ce{NO2}} & \textbf{\ce{CO}} & \textbf{\ce{KD_{10}}} & \textbf{\ce{KD_{2,5}}} & \textbf{Vėjo krypt.} & \textbf{Vėjo greit.} & \textbf{Temp.} \\
                \hline
                2023-08-02    & 06:00:00        & 6,79             & 27,94            & 15,30             & 0,140            & 7,73                  & \textbf{-1,0}          & 240,6                & 6,0                  & 17,7           \\
                2023-08-02    & 07:00:00        & 8,18             & 32,31            & 19,13             & 0,198            & 10,63                 & 1,0                    & 246,1                & 6,0                  & 17,7           \\
                2023-08-02    & 08:00:00        & 8,98             & 45,15            & 37,10             & 0,221            & 9,66                  & 4,0                    & 247,8                & 6,3                  & 18,2           \\
                2023-08-02    & 09:00:00        & 16,56            & 43,03            & 41,12             & 0,221            & 7,73                  & 4,0                    & 241,9                & 6,6                  & 19,0           \\
                2023-08-02    & 10:00:00        & 17,56            & 38,79            & 39,78             & 0,244            & 8,70                  & 3,0                    & 234,7                & 6,9                  & 18,9           \\
                2023-08-02    & 11:00:00        & \textbf{NaN}     & \textbf{NaN}     & \textbf{NaN}      & \textbf{NaN}     & 9,66                  & 6,0           & 225,6 & 5,5 & 16,5 \\
                2023-08-02    & 12:00:00        & \textbf{NaN}     & \textbf{NaN}     & \textbf{NaN}      & \textbf{NaN}     & \textbf{NaN} & \textbf{NaN} & 222,1 & 4,7 & 16,6 \\
                2023-08-02    & 13:00:00        & 10,58            & 50,52            & 51,26             & 0,256            & \textbf{NaN}          & \textbf{NaN}           & 223,2                & 4,8                  & 16,8           \\
                2023-08-02    & 14:00:00        & 5,79             & 72,59            & 59,86             & 0,361            & 24,16                 & 24,0                   & 220,8                & 5,2                  & 16,8           \\
                2023-08-02    & 15:00:00        & 6,98             & 51,01            & 54,51             & 0,303            & 7,73                  & 7,0                    & 204,2                & 4,0                  & 17,9           \\
                2023-08-02    & 16:00:00        & 6,79             & 50,89            & 57,38             & 0,349            & 9,66                  & 11,0                   & 232,1                & 3,8                  & 19,0           \\
                2023-08-02    & 17:00:00        & 16,76            & 16,96            & 37,29             & 0,210            & 6,76                  & 6,0                    & 260,0                & 2,7                  & 18,5           \\
                2023-08-02    & 18:00:00        & 13,57            & 14,34            & 35,19             & 0,221            & 8,70                  & 10,0                   & 313,8                & 3,2                  & 18,0           \\
                \hline
            \end{tabular}
        }
        \label{tab:missing_data}
    \end{table}

    Lentelėje \ref{tab:normalized_missing_data} pateikiama šios iškarpos išvalyta, normalizuota ir transformuota versija.
    6 valanda buvo pašalinta dėl neigiamos \ce{KD_{2,5}} reikšmės, o valandos 11, 12, 13 -- dėl trūkstamų reikšmių.
    \begin{table}[H]
        \centering
        \caption{Normalizuotų duomenų iškarpa su trūkstamomis ir neigiamomis reikšmėmis}
        \small
        \resizebox{\textwidth}{!}{
            \begin{tabular}{|l|c|r|r|r|r|r|r|r|c|r|r|}
                \hline
                \textbf{Data} & \textbf{Val.} & \textbf{\ce{O3}} & \textbf{\ce{NO}} & \textbf{\ce{NO2}} & \textbf{\ce{KD_{10}}} & \textbf{\ce{KD_{2,5}}} & \textbf{Vėjo gr.} & \textbf{Temp.} & \textbf{Dienos kat.} & \textbf{Vėjo $\sin$} & \textbf{Vėjo $\cos$} \\
                \hline
                2023-08-02    & 5             & -0,82            & -0,96            & -0,89             & -0,91                 & -0,95                  & -0,27             & 0,28           & 0                    & -0,88                & -0,47                \\
                2023-08-02    & 7             & -0,87            & -0,93            & -0,88             & -0,89                 & -0,98                  & -0,12             & 0,28           & 0                    & -0,91                & -0,41                \\
                2023-08-02    & 8             & -0,85            & -0,90            & -0,77             & -0,90                 & -0,94                  & -0,08             & 0,31           & 0                    & -0,93                & -0,38                \\
                2023-08-02    & 9             & -0,73            & -0,90            & -0,74             & -0,92                 & -0,94                  & -0,04             & 0,34           & 0                    & -0,88                & -0,47                \\
                2023-08-02    & 10            & -0,71            & -0,91            & -0,75             & -0,91                 & -0,95                  & 0,01              & 0,34           & 0                    & -0,82                & -0,58                \\
                2023-08-02    & 14            & -0,91            & -0,83            & -0,63             & -0,74                 & -0,62                  & -0,24             & 0,24           & 0                    & -0,65                & -0,76                \\
                2023-08-02    & 15            & -0,89            & -0,88            & -0,66             & -0,92                 & -0,89                  & -0,42             & 0,29           & 0                    & -0,41                & -0,91                \\
                2023-08-02    & 16            & -0,89            & -0,88            & -0,64             & -0,90                 & -0,83                  & -0,45             & 0,34           & 0                    & -0,79                & -0,61                \\
                2023-08-02    & 17            & -0,73            & -0,96            & -0,77             & -0,93                 & -0,91                  & -0,61             & 0,32           & 0                    & -0,98                & -0,17                \\
                2023-08-02    & 18            & -0,78            & -0,97            & -0,78             & -0,91                 & -0,84                  & -0,53             & 0,30           & 0                    & -0,72                & 0,69                 \\
                \hline
            \end{tabular}
        }
        \label{tab:normalized_missing_data}
    \end{table}

    \subsection{Valandų blokų filtravimas}
    KNN modelio prognozavimui bus naudojami nepertraukiami penkių valandų įvesties ir dvylikos valandų prognozės blokai.
    Bus taikomas duomenų filtravimo kriterijus: jei pasirinktame laiko bloke (tiek įvesties, tiek prognozės lange) aptinkama bent viena trūkstama reikšmė, visas blokas yra pašalinamas iš tyrimo ir laikomas netinkamu analizei.

    \printbibliography[title = {Literatūra ir šaltiniai}]

\end{document}
